\section{Ausblick}

Aus den Ergebnissen ergeben sich drei naheliegende nächste Schritte zur Fertigstellung des physischen Gesamtsystems MIRA:

\begin{enumerate}
    \item \textbf{Mechatronischer Aufbau des Sortierarms:} Das physische System wird als 3-DOF-Roboterarm (Degrees of Freedom) realisiert. Ein Raspberry Pi führt die Bildverarbeitung aus und sendet die errechneten Objektkoordinaten via Serial-Protokoll an einen ESP32-S3-Mikrocontroller, welcher die Servomotoren (z.B.\ MG996R) präzise ansteuert.
    \item \textbf{Sensorfusion zur Tiefenerfassung:} Da eine monokulare 2D-Kamera keine Tiefeninformationen (Z-Achse) liefern kann, wird ein extrem günstiger Time-of-Flight (ToF) Lasersensor (VL53L0X) direkt am Greifer montiert. Das Kamerasystem steuert den Arm über die X- und Y-Achse direkt über das Objekt, während der Lasersensor die verbleibende Distanz auf der Z-Achse misst, um den Greifvorgang zu triggern.
    \item \textbf{Cloud-basiertes IoT-Dashboard:} Das bereits funktionierende lokale Streamlit-Dashboard wird um Cloud-IoT-Telemetriefunktionen erweitert. Zukünftige Arbeiten konzentrieren sich auf die Implementierung einer verteilten MQTT-Schnittstelle zur asynchronen Übertragung von Sortierstatistiken, historischen Durchsatzraten und Hardware-Diagnosedaten an ein zentrales Monitoring-System zur Fernüberwachung mehrerer dezentraler Sortiereinheiten.
\end{enumerate}